\ParallelTitle{Two languages, one readable page}{两种语言，一份易读文档}

\ParallelSection{choose-route}{Choose a source route}{选择源文件形式}

\ParallelParagraph{opening}{Bilingual PDF places corresponding content in fixed physical columns. The calling agent supplies the text pairs and the order of the document; the renderer aligns their visible starts. A short translation does not have to be padded to match a longer one. Paragraphs, list items and table rows each have their own alignment boundary. Native LaTeX and structured JSON reach the same package, so changing the source route does not create a different design engine. Work in a new project and keep the original input unchanged. This guide is itself an example of the layout it describes.}{Bilingual PDF 将对应内容放在固定的左右栏中。调用模型提供成对的文本及文档顺序，排版工具负责对齐它们在页面上的起始位置。较短的译文不需要为了凑齐长度而补充文字。段落、列表项和表格行各有自己的对齐边界。原生 LaTeX 与结构化 JSON 使用同一个排版包，因此更换输入形式不会产生另一套设计引擎。请在新项目中工作，并保留原始输入。本指南本身就是其所介绍排版方式的示例。}

\begin{ParallelKeep}

\ParallelText{routes-table.row-0}{\phantomsection\label{routes-table}\hypertarget{routes-table}{}\begin{tabularx}{\linewidth}{@{}>{\RaggedRight\arraybackslash}X>{\RaggedRight\arraybackslash}X>{\RaggedRight\arraybackslash}X@{}}\toprule \textbf{Route} & \textbf{Use it when} & \textbf{Result}\\ \midrule \end{tabularx}}{\ifdefstring{\ParallelMode}{right}{\phantomsection\label{routes-table}\hypertarget{routes-table}{}}{}\begin{tabularx}{\linewidth}{@{}>{\RaggedRight\arraybackslash}X>{\RaggedRight\arraybackslash}X>{\RaggedRight\arraybackslash}X@{}}\toprule \textbf{形式} & \textbf{适用情形} & \textbf{结果}\\ \midrule \end{tabularx}}

\ParallelText{routes-table.row-1}{\begin{tabularx}{\linewidth}{@{}>{\RaggedRight\arraybackslash}X>{\RaggedRight\arraybackslash}X>{\RaggedRight\arraybackslash}X@{}}Native LaTeX & You want editable commands, mathematics or standard package hooks. & Compile the manuscript directly.\\ \end{tabularx}}{\begin{tabularx}{\linewidth}{@{}>{\RaggedRight\arraybackslash}X>{\RaggedRight\arraybackslash}X>{\RaggedRight\arraybackslash}X@{}}原生 LaTeX & 需要直接编辑命令、数学公式或标准包钩子。 & 直接编译文稿。\\ \end{tabularx}}

\ParallelText{routes-table.row-2}{\begin{tabularx}{\linewidth}{@{}>{\RaggedRight\arraybackslash}X>{\RaggedRight\arraybackslash}X>{\RaggedRight\arraybackslash}X@{}}Structured JSON & You already have ordered pairs of text, rows and image paths. & Export editable LaTeX, or render it immediately.\\ \end{tabularx}}{\begin{tabularx}{\linewidth}{@{}>{\RaggedRight\arraybackslash}X>{\RaggedRight\arraybackslash}X>{\RaggedRight\arraybackslash}X@{}}结构化 JSON & 已有按顺序组织的文本对、表格行和图像路径。 & 导出可编辑的 LaTeX，或立即渲染。\\ \end{tabularx}}

\ParallelText{routes-table.row-3}{\begin{tabularx}{\linewidth}{@{}>{\RaggedRight\arraybackslash}X>{\RaggedRight\arraybackslash}X>{\RaggedRight\arraybackslash}X@{}}Selected language & The reader needs only the physical left or right side. & Choose left or right mode without rewriting the source.\\ \bottomrule \end{tabularx}}{\begin{tabularx}{\linewidth}{@{}>{\RaggedRight\arraybackslash}X>{\RaggedRight\arraybackslash}X>{\RaggedRight\arraybackslash}X@{}}单语输出 & 读者只需要物理左栏或右栏的内容。 & 选择 left 或 right 模式，无需重写源文件。\\ \bottomrule \end{tabularx}}

\ParallelText{routes-table.caption}{Two source routes, one rendering implementation.}{两种源文件形式，共用一套排版实现。}\end{ParallelKeep}

\ParallelSection{page-flow}{Choose whether paragraphs can cross pages}{选择是否允许段落跨页}

\ParallelParagraph{flow-options}{Set paragraph-flow=breakable in the native preamble to allow ordinary ParallelParagraph content to continue over a page boundary. A paragraph can override that choice with flow=keep. JSON uses layout.paragraph\_flow and a paragraph-level flow field. Keep-together content stays on one page; an oversized unit produces an error rather than being truncated or shrunk. ParallelText and ParallelProse remain explicit bounded and flowing alternatives. The following deliberately continuous paragraph demonstrates the flowing route.}{在原生文稿的导言区设置 paragraph-flow=breakable，即可允许普通 ParallelParagraph 内容跨越分页边界。单个段落可以用 flow=keep 覆盖该设置。JSON 使用 layout.paragraph\_flow 和段落级 flow 字段。保持完整的内容单元会留在一页内；若单元过大，工具会报错，不会截断或缩小它。ParallelText 和 ParallelProse 仍分别提供明确的整块和连续排版方式。下面这个有意保持连续的段落展示跨页排版。}

\ParallelParagraph[flow=breakable]{continuing-prose}{A readable parallel document does not require equal line counts. One language may need more words, different punctuation or a different number of lines to express the supplied content. In a flowing paragraph, each physical column continues at its own natural pace, while the package records the shared paragraph boundary and restores alignment before the next pair begins. This means that a reader can follow a long explanation without a page-sized empty hole caused by moving the whole paragraph to the next sheet. It also means that the last lines of the two translations need not sit at exactly the same height. The useful promise is a stable place to begin the next corresponding unit. When choosing the policy, distinguish continuous prose from objects whose pieces should remain together. A short definition, a formula with a short explanation or a compact diagram may be a good bounded unit. A long explanation that would fill several pages is usually a candidate for flowing prose. Do not put a flowing paragraph inside a minipage or another keep-together wrapper, because the enclosing box prevents page breaking. Do not repair a layout problem by removing sentences, inventing extra text, changing the meaning or shrinking a single troublesome block. Instead, use the documented flow option and inspect the actual result. After building, look at the end of the first page, the continuation on the next page and the start of the following paired section. Check that neither column is clipped, that the physical divider follows the live text block and that the next section begins at a corresponding position. Repeat this check with the requested fonts and language direction. A successful compiler exit proves that a file was produced; it does not prove that all visible page details are correct. Keeping source, language configuration and presentation separate makes a later font, paper or binding change easier to test without rewriting the manuscript.}{可读的双语对照文档并不要求两栏行数相同。不同语言可能需要更多词语、不同的标点，或者不同数量的行，才能表达所提供的内容。在允许跨页的段落中，每一物理栏都按自身自然长度继续排版，排版包记录共同的段落边界，并在下一对内容开始前恢复对齐。因此，读者可以连续阅读长篇说明，而不会因为整段被搬到下一张纸上而留下大片空白。两种语言的最后几行也不必落在完全相同的高度。真正有用的保证，是下一组对应内容有一个稳定的共同起点。选择策略时，应区分连续正文与需要保持完整的对象。简短定义、带简短说明的公式或紧凑图示，通常适合作为不可拆分的单元；可能占据多页的长篇说明，则通常适合连续正文。不要把可跨页段落放进 minipage 或其他保持整块的容器中，因为外层盒子会阻止分页。也不要通过删除句子、凭空补充文字、改变含义或单独缩小某个难排的区块来修补版面。应使用文档规定的跨页选项，并检查实际结果。构建后，查看前一页末尾、下一页的续文，以及其后成对小节的起点。确认两栏都没有被裁切，物理分隔线随实际正文区域延伸，下一节也从对应位置开始。对所需字体和语言方向重复这些检查。编译器成功退出，只能证明生成了文件，并不能证明页面上的所有细节都正确。将源内容、语言配置与呈现设置分开，能够在以后更换字体、纸张或装订方式时，更容易测试而不必重写文稿。}

\ParallelSection{figures-and-references}{Use figures, tables and references}{使用插图、表格与引用}

\ParallelWideFigure{shared-layout}{layout-anatomy.png}{One shared diagram spans the full text width; its captions remain paired.}{同一张示意图横跨整个正文宽度，图注仍然成对对齐。}

\ParallelFigure{localized-pipeline}{pipeline-en.png}{The same pipeline has labels in each column's own language.}{同一处理流程在每一栏中使用该栏语言的标签。}[pipeline-zh-Hans.png]

\ParallelText{quotation}{\begin{quote}“Corresponding starts matter more than identical line counts.” This short quotation is original example text; quotation layout does not imply an external source.\end{quote}}{\begin{quote}“对应的起点比相同的行数更重要。”这段短引文是原创示例文本；引文样式不表示它来自外部来源。\end{quote}}

\begin{ParallelKeep}\ParallelText{column-width}{For a usable text width W and a gap g, each paired column has width w. If W is 170 mm and g is 6 mm, w is 82 mm.}{设可用正文宽度为 W，栏间距为 g，则每一对照栏的宽度为 w。如果 W 为 170 mm，g 为 6 mm，则 w 为 82 mm。}\ParallelText{column-width.formula}{\[ w=\frac{W-g}{2} \]}{\[ w=\frac{W-g}{2} \]}\end{ParallelKeep}

\ParallelSection{check-and-deliver}{Check the PDF and deliver the project}{检查 PDF 并交付项目}

\ParallelText{checks.item-1}{\begin{itemize}\item \phantomsection\label{checks}\hypertarget{checks}{}Inspect actual pages for matching starts, table-row alignment, paragraph continuations, clipping and intended whitespace. A first-page preview cannot show the whole document.\end{itemize}}{\begin{itemize}\item \ifdefstring{\ParallelMode}{right}{\phantomsection\label{checks}\hypertarget{checks}{}}{}查看实际页面，检查对应起点、表格行对齐、段落续排、裁切及预期留白。首页预览无法展示整份文档。\end{itemize}}

\ParallelText{checks.item-2}{\begin{itemize}\item For RTL and CJK, inspect shaping, punctuation, mixed-script identifiers and physical column order. Language direction does not decide which side a language occupies.\end{itemize}}{\begin{itemize}\item 对于 RTL 与 CJK，检查字形连接、标点、混合文字标识符和物理栏顺序。语言书写方向并不决定它位于哪一栏。\end{itemize}}

\ParallelText{checks.item-3}{\begin{itemize}\item Follow generated references to their real destinations. Printed folios can differ from physical PDF page numbers. Keep relative cross-document targets beside each other.\end{itemize}}{\begin{itemize}\item 沿自动生成的引用检查真实目标。印刷页码可能不同于 PDF 的物理页号。使用相对跨文档链接的文件应保存在一起。\end{itemize}}

\ParallelText{checks.item-4}{\begin{itemize}\item Deliver the PDF and a portable source project with required package, image and license files. Native TeX executes code; disabled shell escape is not a filesystem sandbox.\end{itemize}}{\begin{itemize}\item 交付 PDF 及可移植源项目，包括必要的排版包、图像与许可证文件。原生 TeX 可以执行代码；关闭 shell escape 并不等于文件系统沙箱。\end{itemize}}

\ParallelText{back-to-flow}{\ParallelReference{page-flow}{Return to the page-breaking controls}}{\ParallelReference{page-flow}{返回段落分页设置}}
